\RequirePackage[l2tabu, orthodox]{nag}
\RequirePackage{silence}
\WarningFilter{fmtcount}{\ordinal already defined use \FCordinal instead}
\documentclass[english, french]{beamer}
%INSTALL

%avoids a warning
\usepackage[log-declarations=false]{xparse}
\usepackage{fontspec} %font selecting commands
\usepackage{xunicode}
%warn about missing characters
\tracinglostchars=2

%REDAC
\usepackage{booktabs}
\usepackage{calc}

\usepackage{mathtools} %load this before babel!
	\mathtoolsset{showonlyrefs,showmanualtags}

\usepackage{babel}
%suppresses the warning about frenchb not modifying the captions (“—” to “:” in “Figure 1 – Legend”).
	\frenchbsetup{AutoSpacePunctuation=false,SuppressWarning=false}

\usepackage[super]{nth}
\usepackage{listings} %typeset source code listings
	\lstset{language=XML,tabsize=2,literate={"}{{\tt"}}1,captionpos=b}
\usepackage[nolist,smaller,printonlyused]{acronym}%,smaller option produces warnings from relsize in some cases, it seems.
\usepackage[nodayofweek]{datetime}%must be loaded after the babel package
\usepackage{xspace}
\usepackage{hyperref}% option pdfusetitle must be introduced here, not in hypersetup.
%breaklinks makes links on multiple lines into different PDF links to the same target.
%colorlinks (false): Colors the text of links and anchors. The colors chosen depend on the the type of link. In spite of colored boxes, the colored text remains when printing.
%linkcolor=black: this leaves other links in colors, e.g. refs in green, don't print well.
%pdfborder (0 0 1, set to 0 0 0 if colorlinks): width of PDF link border
%hidelinks
\hypersetup{breaklinks,bookmarksopen,colorlinks=true,urlcolor=blue,linkcolor=,hyperfigures=true}
% hyperref doc says: Package bookmark replaces hyperref’s bookmark organization by a new algorithm (...) Therefore I recommend using this package.
\usepackage{bookmark}

% center floats by default, but do not use with float
% \usepackage{floatrow}
% \makeatletter
% \g@addto@macro\@floatboxreset\centering
% \makeatother
\usepackage{ragged2e} %new com­mands \Cen­ter­ing, \RaggedLeft, and \RaggedRight and new en­vi­ron­ments Cen­ter, FlushLeft, and FlushRight, which set ragged text and are eas­ily con­fig­urable to al­low hy­phen­ation (the cor­re­spond­ing com­mands in LaTeX, all of whose names are lower-case, pre­vent hy­phen­ation al­to­gether). 
\usepackage{siunitx} %[expproduct=tighttimes, decimalsymbol=comma]
\sisetup{detect-all}% to detect e.g. when in math mode (use a math font)
\usepackage{braket} %for \Set
\usepackage{natbib}

\usepackage{amsmath,amsthm}
% \usepackage{amsfonts} %not required?
% \usepackage{dsfont} %for what?
%unicode-math overwrites the following commands from the mathtools package: \dblcolon, \coloneqq, \Coloneqq, \eqqcolon. Using the other colon-like commands from mathtools will lead to inconsistencies. Plus, Using \overbracket and \underbracke from mathtools package. Use \Uoverbracket and \Uunderbracke for original unicode-math definition.
%use exclusively \mathbf and choose math bold style below.
\usepackage[warnings-off={mathtools-colon, mathtools-overbracket}, bold-style=ISO]{unicode-math}

\defaultfontfeatures{
	Fractions=On,
	Mapping=tex-text% to turn "--" into dashes, useful for bibtex%%
}
\defaultfontfeatures[\rmfamily, \sffamily]{
	Fractions=On,
	Mapping=% to leave " alone (disable the default mapping tex-text; requires loading the font afterwards?)
}
\newfontfamily\xitsfamily{XITS}
\newfontfamily\texgyretermesfamily{TeX Gyre Termes}
\newfontfamily\texgyreherosfamily{TeX Gyre Heros}
\newfontfamily\lmfamily{Latin Modern Roman}
\setmainfont{Latin Modern Roman}
\setsansfont{Latin Modern Sans}
\setmonofont{Latin Modern Mono}
\defaultfontfeatures{}%disable default font features to avoid warnings with math fonts.
\setmathfont{XITS Math}
\setmathfont[range={\mathcal,\mathbfcal},StylisticSet=1]{XITS Math}

\usepackage{cleveref}% cleveref should go "laster" than hyperref
%GRAPHICS
\usepackage{pgf}
\usepackage{pgfplots}
	\usetikzlibrary{matrix,fit,plotmarks,calc,trees,shapes.geometric,positioning,plothandlers}
\pgfplotsset{compat=1.11}
\usepackage{graphicx}

\graphicspath{{graphics/},{graphics-dm/}}
\DeclareGraphicsExtensions{.pdf}

%HACKING
\usepackage{printlen}
\uselengthunit{mm}
% 	\newlength{\templ}% or LenTemp?
% 	\setlength{\templ}{6 pt}
% 	\printlength{\templ}
\usepackage{etoolbox} %for addtocmd command
\usepackage{scrhack}% load at end. Corrects a bug in float package, which is outdated but might be used by other packages
\usepackage{xltxtra} %somebody said that this is loaded by fontspec, but does not seem correct: if not loaded explicitly, does not appear in the log and \showhyphens is not corrected.

%Beamer-specific
%ADD
\usepackage{appendixnumberbeamer}
%\setbeamersize{text margin left=0.1cm, text margin right=0.1cm} 
\setbeamertemplate{navigation symbols}{}
\usetheme{BrusselsBelgium}
\usefonttheme{professionalfonts}


\newcommand{\R}{ℝ}
\newcommand{\N}{ℕ}
\newcommand{\Z}{ℤ}
\newcommand{\card}[1]{\lvert{#1}\rvert}
\newcommand{\powerset}[1]{\mathscr{P}(#1)}%\mathscr rather than \mathcal: scr is rounder than cal (at least in XITS Math).
\newcommand{\suchthat}{\;\ifnum\currentgrouptype=16 \middle\fi|\;}
%\newcommand{\Rplus}{\reels^+\xspace}

\AtBeginDocument{%
	\renewcommand{\epsilon}{\varepsilon}
% we want straight form of \phi for mathematics, as recommended in UTR #25: Unicode support for mathematics.
%	\renewcommand{\phi}{\varphi}
}

% with amssymb, but I don’t want to use amssymb just for that.
% \newcommand{\restr}[2]{{#1}_{\restriction #2}}
%\newcommand{\restr}[2]{{#1\upharpoonright}_{#2}}
\newcommand{\restr}[2]{{#1|}_{#2}}%sometimes typed out incorrectly within \set.
%\newcommand{\restr}[2]{{#1}_{\vert #2}}%\vert errors when used within \Set and is typed out incorrectly within \set.
\DeclareMathOperator*{\argmax}{arg\,max}
\DeclareMathOperator*{\argmin}{arg\,min}


%ARG TH
\newcommand{\AF}{\mathcal{AF}}
\newcommand{\labelling}{\mathcal{L}}
\newcommand{\labin}{\textbf{in}\xspace}
\newcommand{\labout}{\textbf{out}}
\newcommand{\labund}{\textbf{undec}\xspace}
\newcommand{\nonemptyor}[2]{\ifthenelse{\equal{#1}{}}{#2}{#1}}
\newcommand{\gextlab}[2][]{
	\labelling{\mathcal{GE}}_{(#2, \nonemptyor{#1}{\ibeatsr{#2}})}
}
\newcommand{\allargs}{A^*}
\newcommand{\args}{A}
\newcommand{\ar}{a}

%MCDA+Arg
\newcommand{\dm}{d}
\newcommand{\ileadsto}{\leadsto}
\newcommand{\mleadsto}[1][\eta]{\leadsto_{#1}}
\newcommand{\ibeats}{\vartriangleright}
\newcommand{\mbeats}[1][\eta]{\vartriangleright_{#1}}


%MISC
\newcommand{\lequiv}{\Vvdash}
\newcommand{\weightst}{W^{\,t}}

%MCDA classical
\newcommand{\crits}{\mathcal{J}}
\newcommand{\altspace}{\mathbb{A}}
\newcommand{\alts}{A}

%Sorting
\newcommand{\cats}{\mathcal{C}}
\newcommand{\catssubsets}{2^\cats}
\newcommand{\catgg}{\vartriangleright}
\newcommand{\catll}{\vartriangleleft}
\newcommand{\catleq}{\trianglelefteq}
\newcommand{\catgeq}{\trianglerighteq}
\newcommand{\alttoc}[2][x]{(#1 \xrightarrow{} #2)}
\newcommand{\alttocat}[3]{(#2 \xrightarrow{#1} #3)}
\newcommand{\alttoI}{(x \xrightarrow{} \left[\underline{C_x}, \overline{C_x}\right])}
\newcommand{\alttocatdm}[3][t]{\left(#2 \thinspace \raisebox{-3pt}{$\xrightarrow{#1}$}\thinspace #3\right)}
\newcommand{\alttocatatleast}[2]{\left(#1 \thinspace \raisebox{-3pt}{$\xrightarrow[]{≥}$}\thinspace #2\right)}
\newcommand{\alttocatatmost}[2]{\left(#1 \thinspace \raisebox{-3pt}{$\xrightarrow[]{≤}$}\thinspace #2\right)}

\newcommand{\source}{\scriptsize}
\newcommand{\commentOC}[1]{{\selectlanguage{french}{\todo{OC : #1}}}}
%Or: \todo[color=green!40]

%this probably requires outdated float package, see doc KomaScript for an alternative.
% \newfloat{program}{t}{lop}
% \floatname{program}{PM}

%\crefname{axiom}{axiom}{axioms}%might be needed for workaround bug in cref when defining new theorems?

%\ifdefined\theorem\else
%\newtheorem{theorem}{\iflanguage{english}{Theorem}{Théorème}}
%\fi

%which line breaks are chosen: accept worse lines, therefore reducing risk of overfull lines. Default = 200
\tolerance=2000
%accept overfull hbox up to...
\hfuzz=2cm
%reduces verbosity about the bad line breaks
\hbadness 5000
%sloppy sets tolerance to 9999
\apptocmd{\sloppy}{\hbadness 10000\relax}{}{}

% WRITING
%\newcommand{\ie}{i.e.\@\xspace}%to try
%\newcommand{\eg}{e.g.\@\xspace}
%\newcommand{\etal}{et al.\@\xspace}
\newcommand{\ie}{i.e.\ }
\newcommand{\eg}{e.g.\ }
\newcommand{\mkkOK}{\checkmark}%\color{green}{\checkmark}
\newcommand{\mkkREQ}{\ding{53}}%requires pifont?%\color{green}{\checkmark}
\newcommand{\mkkNO}{}%\text{\color{red}{\textsf{X}}}

\makeatletter
\newcommand{\boldor}[2]{%
	\ifnum\strcmp{\f@series}{bx}=\z@
		#1%
	\else
		#2%
	\fi
}
\newcommand{\textstyleElProm}[1]{\boldor{\MakeUppercase{#1}}{\textsc{#1}}}
\makeatother
\newcommand{\electre}{\textstyleElProm{Électre}\xspace}
\newcommand{\electreIv}{\textstyleElProm{Électre Iv}\xspace}
\newcommand{\electreIV}{\textstyleElProm{Électre IV}\xspace}
\newcommand{\electreIII}{\textstyleElProm{Électre III}\xspace}
\newcommand{\electreTRI}{\textstyleElProm{Électre Tri}\xspace}
% \newcommand{\utadis}{\texorpdfstring{\textstyleElProm{utadis}\xspace}{UTADIS}}
% \newcommand{\utadisI}{\texorpdfstring{\textstyleElProm{utadis i}\xspace}{UTADIS I}}

%TODO
% \newcommand{\textstyleElProm}[1]{{\rmfamily\textsc{#1}}} 


\newlength{\GraphsNodeSep}
\setlength{\GraphsNodeSep}{7mm}

% MCDA Drawing Sorting
\newlength{\MCDSCatHeight}
\setlength{\MCDSCatHeight}{6mm}
\newlength{\MCDSAltHeight}
\setlength{\MCDSAltHeight}{4mm}
%separation between two vertical alts
\newlength{\MCDSAltSep}
\setlength{\MCDSAltSep}{2mm}
\newlength{\MCDSCatWidth}
\setlength{\MCDSCatWidth}{3cm}
\newlength{\MCDSEvalRowHeight}
\setlength{\MCDSEvalRowHeight}{6mm}
\newlength{\MCDSAltsToCatsSep}
\setlength{\MCDSAltsToCatsSep}{1.5cm}
\newcounter{MCDSNbAlts}
\newcounter{MCDSNbCats}
\newlength{\MCDSArrowDownOffset}
\setlength{\MCDSArrowDownOffset}{0mm}

\tikzset{/Graphs/dot/.style={
	shape=circle, fill=black, inner sep=0, minimum size=1mm
}}
\tikzset{/MC/D/S/alt/.style={
	shape=rectangle, draw=black, inner sep=0, minimum height=\MCDSAltHeight, minimum width=2.5cm, anchor=north east
}}
\tikzset{MC/D/S/pref/.style={
	shape=ellipse, draw=gray, thick
}}
\tikzset{/MC/D/S/cat/.style={
	shape=rectangle, draw=black, inner sep=0, minimum height=\MCDSCatHeight, minimum width=\MCDSCatWidth, anchor=north west
}}
\tikzset{/MC/D/S/evals matrix/.style={
	matrix, row sep=-\pgflinewidth, column sep=-\pgflinewidth, nodes={shape=rectangle, draw=black, inner sep=0mm, text depth=0.5ex, text height=1em, minimum height=\MCDSEvalRowHeight, minimum width=12mm}, nodes in empty cells, matrix of nodes, inner sep=0mm, outer sep=0mm, row 1/.style={nodes={draw=none, minimum height=0em, text height=, inner ysep=1mm}}
}}

% Beliefs
\tikzset{/Beliefs/D/S/attacker/.style={
	shape=rectangle, draw, minimum size=8mm
}}
\tikzset{/Beliefs/D/S/supporter/.style={
	shape=circle, draw
}}

\newcommand{\tikzmark}[1]{%
	\tikz[overlay, remember picture, baseline=(#1.base)] \node (#1) {};%
}


\begin{acronym}
\acro{AMCD}{Aide Multicritère à la Décision}
\acro{ASA}{Argument Strength Assessment}
\acro{DA}{Decision Analysis}
\acro{DM}{Decision Maker}
\acro{DPr}{Deliberated Preferences}
\acro{DRSA}{Dominance-based Rough Set Approach}
\acro{DSS}{Decision Support Systems}
\acrodefplural{DSS}{Decision Support Systems}
% \newacroplural{DSS}[DSSes]{Decision Support Systems}
\acro{EJOR}{European Journal of Operational Research}
\acro{LNCS}{Lecture Notes in Computer Science}
\acro{MCDA}{Multicriteria Decision Aid}
\acro{MIP}{Mixed Integer Program}
\acro{NCSM}{Non Compensatory Sorting Model}
\acro{PL}{Programme Linéaire}
\acro{PLNE}{Programme Linéaire en Nombres Entiers}
\acro{PM}{Programme Mathématique}
\acro{MP}{Mathematical Program}
\acro{MIP}{Mixed Integer Program}
% \newacroplural{PM}{Programmes Mathématiques}
%acrodefplural since version 1.35, my debian has \ProvidesPackage{acronym}[2009/01/25, v1.34, Support for acronyms (Tobias Oetiker)]
\acrodefplural{PM}{Programmes Mathématiques}
\acro{PMML}{Predictive Model Markup Language}
\acro{RESS}{Reliability Engineering \& System Safety}
\acro{SMAA}{Stochastic Multicriteria Acceptability Analysis}
\acro{URPDM}{Uncertainty and Robustness in Planning and Decision Making}
\acro{XML}{Extensible Markup Language}
\end{acronym}

%this approach does not generalize to multipart nodes.
%\tikzset{/uml/interface/.style={rectangle, draw, fill=yellow!20, align=left, node font=\fontspec{Latin Modern Mono Light}, node contents={$<<$interface$>>$\\\textbf{#1}}, name=#1
%}}
%this approach does not apply for <<interface>>. I’ll do it manually, as I can’t find anything better. Use: \nodepart[font=]{one} \small <<interface>>\\\bfseries ItemDAO
\tikzset{/uml/class/.style={rectangle, draw, align=center, fill=yellow!20, node font=\ttfamily, font=\bfseries
}}
\tikzset{/uml/abstract class/.style={rectangle, draw, align=center, fill=yellow!20, node font=\fontspec{Latin Modern Mono Light}, font=\bfseries\itshape
}}
\tikzset{/uml/class3/.style={
rectangle split, rectangle split every empty part={}, rectangle split parts=3, rectangle split part align={center, left, left}, draw, fill=yellow!20, rectangle split empty part height=0, node font=\ttfamily, every one node part/.style={font=\bfseries, align=center}, every two node part/.style={align=left}, every three node part/.style={align=left}
}}
\tikzset{/uml/extends/.style={draw, -open triangle 45}}
\tikzset{/uml/implements/.style={draw, -open triangle 45, dashed}}
%prefix= used to cancel \texttt, otherwise this cancels the effect of the font selection \fontspec{Latin Modern Mono Light}, and thus renders \bfseries inoperant. With prefix=, the result is correct.
%\jeeref[prefix=]{javax.faces.component.html/HtmlCommandButton}\\
\tikzset{/uml/table/.style={rectangle, draw, align=center, fill=blue!20, node font=\ttfamily, font=\bfseries
}}
\tikzset{/uml/table2/.style={
rectangle split, rectangle split every empty part={}, rectangle split parts=2, rectangle split part align={center, left}, draw, fill=blue!20, rectangle split empty part height=0, node font=\ttfamily, every one node part/.style={font=\bfseries, align=center}, every two node part/.style={align=left}
}}
\tikzset{/uml/dbkey/.style={draw, ->}}
%http://tex.stackexchange.com/questions/79781/placing-anchor-before-and-after-text-in-multipart-rectangle



\newbool{refAPIIsAnnotation}
\newbool{refAPIDisplayFull}
\newbool{refAPIOmitFQName}
\pgfkeys{
	/refAPI/.is family,
	/refAPI/.cd,
	base url/.estore in = \refAPIBaseUrl,
	@/.is if = refAPIIsAnnotation,
	full/.is if = refAPIDisplayFull,
	omitfqname/.is if = refAPIOmitFQName,
	prefix/.store in = \refAPIPrefix,
	suffix/.estore in = \refAPISuffix,
	default/.style = {
		base url = ,
		@ = false,
		full = false,
		omitfqname = false,
		prefix = \texttt,
		suffix = ,
	},
}

%Usage: \refAPIParseRef{javax.persistence/EntityManager} ; \refAPIParseRef{javax.persistence/PersistenceContextType\#EXTENDED}
%Defines: \refAPIFQName (before #) ; \refAPIFieldFull (after #) ; \refAPIFieldShort (after # and before -) ; \refAPIPackageSlashes (FQName before / with slashes instead of dots) ; \refAPIClass (FQName after /)
\newcommand{\refAPIParseRef}[1]{%
	\IfSubStr{#1}{\#}{%
		\StrBefore{#1}{\#}[\refAPIFQName]%
		\StrBehind{#1}{\#}[\refAPIFieldFull]%
	}{%
		\edef\refAPIFQName{#1}%
		\edef\refAPIFieldFull{}%
	}%
	\IfSubStr{\refAPIFieldFull}{-}{%
		\StrBefore{\refAPIFieldFull}{-}[\refAPIFieldShort]%
	}{%
		\edef\refAPIFieldShort{\refAPIFieldFull}%
	}%
	\StrBefore{\refAPIFQName}{/}[\refAPIPackage]%
	\StrBehind{\refAPIFQName}{/}[\refAPIClass]%
	\StrSubstitute{\refAPIPackage}{.}{/}[\refAPIPackageSlashes]%
}

%Usage: \jeeref{javax.persistence/EntityManager} ; \jeeref[@]{javax.persistence/PersistenceContextType\#EXTENDED}
\newcommand{\jeeref}[2][]{\refAPIParseRef{#2}\jrefAPI{base url = https://docs.oracle.com/javaee/7/api/, #1}}
\newcommand{\jseref}[2][]{\refAPIParseRef{#2}\jrefAPI{base url = https://docs.oracle.com/javase/8/docs/api/, #1}}
\newcommand{\jsftag}[2][]{\refAPIParseRef{#2}\jrefAPI{base url = https://docs.oracle.com/javaee/7/javaserver-faces-2-2/vdldocs-facelets/, full, prefix=:, #1}}
%Standard HTML RenderKit : https://docs.oracle.com/javaee/7/javaserver-faces-2-2/renderkitdocs/HTML_BASIC/javax.faces.Inputjavax.faces.Text.html

%Accepts keys: @, prefix, suffix. Requires key: base url. Expects defined: \refAPIField ; \refAPIPackageSlashes ; \refAPIClass
%Usage: \texttt{"Hello".\jseref[omitfqname]{java.lang/String\#substring-int-int-}(1, 3);}
\newcommand{\jrefAPI}[1]{%
	\pgfkeys{/refAPI/.cd, default, #1}%
	\ifbool{refAPIIsAnnotation}{%
		\edef\refAPIAnnot{@}%
	}{%
		\edef\refAPIAnnot{}%
	}%
	\IfEq{\refAPIFieldFull}{}{%
		{%
			\ifbool{refAPIDisplayFull}{%
				\refAPIPackage%
			}{%
			}%
			\refAPIPrefix{%
			\href{%
				\refAPIBaseUrl\refAPIPackageSlashes/\refAPIClass.html%
			}{%
				\refAPIAnnot\refAPIClass%
			}%
		}\refAPISuffix}%
	}{%
		{\refAPIPrefix{%
			\ifbool{refAPIOmitFQName}{%
			}{%
				\refAPIAnnot\refAPIClass.%
			}%
			\href{%
				\refAPIBaseUrl\refAPIPackageSlashes/\refAPIClass.html\#\refAPIFieldFull%
			}{%
				\refAPIFieldShort%
			}%
		}\refAPISuffix}%
	}%
}



\setbeamertemplate{headline}[singleline]
\setbeamertemplate{footline}[onlypage]

\title{JDBC}
\subject{Databases}
\keywords{Transaction}
\author{Olivier Cailloux}
\institute[LAMSADE]{LAMSADE, Université Paris-Dauphine}
\date{Version du \today}

\begin{document}
\bibliographystyle{apalike}

\begin{frame}[plain]
	\tikz[remember picture,overlay]{
		\path (current page.south west) node[anchor=south west, inner sep=0] {
			\includegraphics[height=1cm]{LAMSADE95.jpg}
		};
		\path (current page.south) ++ (0, 1mm) node[anchor=south, inner sep=0] {
			\includegraphics[height=9mm]{Dauphine.jpg}
		};
		\path (current page.south east) node[anchor=south east, inner sep=0] {
			\includegraphics[height=1cm]{PSL.png}
		};
	}
   \titlepage
\end{frame}
\addtocounter{framenumber}{-1}

\section{Présentation}
\begin{frame}
	\frametitle{Présentation}
	\begin{itemize}
		\item JDBC ? \pause Java Database Connectivity \pause
		\item Une API pour se connecter à des données relationnelles
	\end{itemize}
	\begin{minipage}[t]{4cm}
		\begin{itemize}
			\item Programmation indépendante du fournisseur de BD
			\item App. programmée via API JDBC
			\item App. inclut pilotes du fournisseur
			\item Ces pilotes font la traduction
		\end{itemize}
	\end{minipage}\hfill%
	\begin{minipage}[t]{6.6cm}
		\begin{tikzpicture}[baseline=(App.north)]
			\path node[draw, ellipse] (App) {Application};
			\path (App.south) ++(0, -3mm) node[draw, ellipse, anchor=north] (JDBC) {API JDBC};
			\path (JDBC.south) ++(-1.4cm, -3mm) node[draw, ellipse, anchor=north] (Driver1) {Pilote Pg.};
			\path (Driver1.south) ++(0, -6mm) node[draw, ellipse, anchor=north] (BD1) {BD Postgresql};
			\path (JDBC.south) ++(1.4cm, -3mm) node[draw, ellipse, anchor=north] (Driver2) {Pilote H2};
			\path (Driver2.south) ++(0, -2mm) node[draw, ellipse, anchor=north] (BD2) {BD en mémoire};
			\path[<->, draw] (App) -- (JDBC);
			\path[<->, draw] (JDBC) -- (Driver1);
			\path[<->, draw] (JDBC) -- (Driver2);
			\path[<->, draw] (Driver1) -- (BD1);
			\path[<->, draw] (Driver2) -- (BD2);
		\end{tikzpicture}
	\end{minipage}
\end{frame}

\section{Instanciation}
\begin{frame}
	\frametitle{Instantiation}
	\begin{itemize}
		\item Souhait : instancier pilote adéquat avec minimum de code spécifique à un fournisseur
		\item API JDBC nous fournit \jseref{java.sql/DriverManager}
		\item Appeler \texttt{DriverManager.getConnection(String url)}
		\item \texttt{url} au format \texttt{jdbc:subprotocol:subname}
		\item Exemple : \texttt{jdbc:postgresql:mydb} (cf. \href{https://jdbc.postgresql.org/documentation/head/connect.html}{Doc JDBC Postgresql})
	\end{itemize}
	Mais comment ça marche ?
\end{frame}

\begin{frame}
	\frametitle{Fonctionnement de l’instantiation}
	\begin{itemize}
		\item Le pilote fournisseur est inclus aux bibliothèques runtime de l’application
		\item Le JAR pilote inclut un fichier nommé (par convention) \texttt{META-INF/services/java.sql.Driver}
		\item Ce fichier nomme la classe que \texttt{DriverManager} doit charger
		\item \texttt{DriverManager} charge toutes ces classes (si plusieurs pilotes accessibles)
		\item Ayant l’URL, \texttt{DriverManager} cherche un pilote enregistré qui peut la lire
		\item Il instancie ce pilote et le renvoie à l’appelant ou l’utilise en arrière-plan
	\end{itemize}
	
\end{frame}

\begin{frame}
	\frametitle{Remarques concernant l’instanciation}
	\begin{itemize}
		\item Avec \texttt{DriverManager} on peut aussi obtenir le \jseref{java.sql/Driver} (utile pour avoir n° de version par exemple)
		\item Beaucoup de tutoriels sur le net suggèrent d’enregistrer explicitement le pilote {\tiny par exemple avec \texttt{Class.forName()}}. Ce n’est plus nécessaire depuis longtemps (cf. explication précédente).
	\end{itemize}
\end{frame}

\section{Modèle relationnel}
\begin{frame}
	\frametitle{Modèle relationnel}
	\begin{itemize}
		\item JDBC permet d’accéder aux BD
		\item Suit le modèle relationnel
		\item Élément central du modèle : la relation (une table, par exemple)
		\item Relations produites par des JOIN, SELECT, etc.
	\end{itemize}
	\begin{block}{Forces du modèle relationnel}
		\pause
		\begin{itemize}
			\item Garanties théoriques (algèbre relationnelle)
			\item Efficace
			\item Standard de fait depuis \textasciitilde 1990
			\item Robuste : \href{https://dl.acm.org/citation.cfm?id=362685}{Codd 1970}, ANSI (puis ISO) SQL \href{https://www.iso.org/iso/home/store/catalogue_ics/catalogue_detail_ics.htm?csnumber=16661}{1987} ; … ; \href{https://www.iso.org/iso/home/store/catalogue_ics/catalogue_detail_ics.htm?csnumber=53681}{2011}
		\end{itemize}
	\end{block}
	Bémol : nombreuses variations propriétaires
\end{frame}

\section{Vue d’ensemble}
\begin{frame}
	\frametitle{Vue d’ensemble de l’API JDBC}
	\begin{itemize}
		\item Fournisseur de SGBD implémente un pilote JDBC
		\item Obtenir une \jseref{java.sql/Connection} (communique avec le pilote)
		\item \texttt{Connection} permet les \emph{transactions}
		\item Transaction : ensemble atomique de \og{}statements\fg{} SQL
		\item Gérer début et fin de transaction via la connexion
		\item Exécuter des statements SQL via cette connexion
		\item Par défaut, mode auto-commit : une transaction par stmt
		\item Via \texttt{Connection} : exécution requêtes, navigation de \href{https://docs.oracle.com/javase/8/docs/api/java/sql/ResultSet.html}{\texttt{ResultSet}}s, …
%		\item (Pilote JDBC peut aussi supporter JTA)
		\item Puis \emph{fermer} la connexion
	\end{itemize}
	Cf. \href{https://docs.oracle.com/javase/tutorial/jdbc/}{tutoriel}
\end{frame}

\section{Instanciation}
\begin{frame}
	\frametitle{Instanciation}
	\begin{block}{Approche 1 (Java SE, typiquement)}
		\begin{itemize}
			\item Une classe fournisseur implémente \href{https://docs.oracle.com/javase/8/docs/api/java/sql/Driver.html}{\texttt{Driver}}
			\item Développeur appelle \jseref{java.sql/DriverManager}
			\item \texttt{DriverManager} trouve le pilote et l’instancie
			\item Exemple : \texttt{DriverManager.getConnection(url)}
		\end{itemize}
	\end{block}
	\begin{block}{Approche 2 (Java EE, typiquement)}
		\begin{itemize}
			\item Une classe fournisseur implémente \jseref{javax.sql/DataSource}
			\item Source accessible via JNDI à un endroit convenu
			\item Développeur instancie DataSource par \texttt{lookup} JNDI puis appelle \texttt{source.getConnection(…)}
		\end{itemize}
	\end{block}
\end{frame}

\begin{frame}
	\frametitle{Injection de la \texttt{DataSource} (Java EE)}
	\begin{itemize}
		\item Injection de ressources via \jeeref[@]{javax.annotation/Resource}
		\item Le conteneur va chercher la ressource via JNDI
		\item Nom JNDI par défaut selon type de la ressource
		\item Pour nous : \texttt{@Resource DataSource myDataSource;}
	\end{itemize}
\end{frame}

\section{Exploitation}
\subsubsection{Commandes et résultats}
\begin{frame}
	\frametitle{\texttt{Statement} et \texttt{ResultSet}}
	\begin{itemize}
		\item Création d’un \jseref{java.sql/Statement} (via \texttt{Connection})
		\item Via \texttt{Statement} : exécution d’une commande SQL (SELECT, UPDATE…)
		\item Via \texttt{Statement} : paramétrisation possible (nb résultats max…)
		\item Obtention (si SELECT) d’un \jseref{java.sql/ResultSet}
		\item \texttt{ResultSet} associé à une ligne courante ; initialement : avant la première
		\item Naviguer via \texttt{next()} aux lignes suivantes
		\item Invoquer \texttt{getInt(columnLabel)}, \texttt{getString(columnLabel)}…
	\end{itemize}
\end{frame}

\begin{frame}[fragile]
	\frametitle{\texttt{PreparedStatement}}
	\begin{itemize}
		\item \jseref{java.sql/PreparedStatement} : précompilé + paramétrisation facile
		\item La commande SQL contient des \og{}\texttt{?}\fg{}
		\item Invoquer \texttt{setInt}, \texttt{setString}… pour les paramètres
	\end{itemize}
	\begin{exampleblock}{Exemple \texttt{PreparedStatement}}
		\begin{lstlisting}[basicstyle=\small\NoAutoSpacing\ttfamily, aboveskip=0pt, belowskip=0pt, showstringspaces=false]
String s="update USER set NAME = ? where ID = ?";
PreparedStatement stmt = con.prepareStatement(s);
stmt.setString(1, "NewName");
stmt.setInt(2, 1234);
boolean isResultSet = stmt.execute();
assert(!isResultSet);
assert(stmt.getUpdateCount() == 1);
		\end{lstlisting}
	\end{exampleblock}
	Utiliser \texttt{PreparedStatement} pour éviter les attaques de type injection SQL !
\end{frame}

\subsubsection{Transactions}
\begin{frame}
	\frametitle{Transactions}
	Par défaut, mode \emph{auto commit} : une transaction par commande
	\begin{block}{Gestion de transactions explicite}
		\begin{itemize}
			\item Invoquer \texttt{setAutoCommit} sur \jseref{java.sql/Connection}
			\item Exécuter les commandes normalement
			\item Puis invoquer \texttt{commit} sur \texttt{Connection}
			\item Ou : \texttt{rollback}
			\item Voir aussi : \texttt{getTransactionIsolation}, \texttt{setTransactionIsolation}
		\end{itemize}
	\end{block}
\end{frame}

\section{Outils}
\begin{frame}
	\frametitle{PostgreSQL}
	\begin{itemize}
		\item Installer PostgreSQL (\href{https://www.postgresql.org/}{site} {\tiny ou \texttt{sudo apt-get install postgresql}})
		\item Possible d’utiliser l’interface graphique d’administration pgAdmin {\tiny voir \href{https://wiki.postgresql.org/wiki/Troubleshooting_Installation\#Collect_the_installer_log_file}{logs} si nécessaire}
		\item Instructions ci-dessous pour ligne de commande linux, adapter pour autres OS
		\item En mode privilégié : \texttt{sudo -u postgres bash}
		\begin{itemize}
			\item Se créer un utilisateur avec mot de passe : \texttt{createuser -P user}
			\item Créer une base de données à laquelle cet utilisateur a accès : \texttt{createdb -O user db}
		\end{itemize}
		\item Test connexion : psql db (ok sans mot de passe)
		\item Test connexion réseau : psql -h localhost db (exige mot de passe)
		\item Droits de connexion : voir \path{/etc/postgresql/9.4/main/} \path{pg_hba.conf} {\tiny changer ligne IPv6 local connections puis reload}
	\end{itemize}
\end{frame}

\begin{frame}
	\frametitle{Serveur d’application}
	\begin{itemize}
		\item Avec un serveur d’application, il faut renseigner le pilote JDBC utilisé dans JNDI
		\item Conseil : s’assurer d’abord que l’instanciation fonctionne avec un projet simple Java SE et Maven
		\item Via l’interface d’administration, indiquer la source à utiliser pour la connexion JNDI par défaut
		\item Inclure le pilote JDBC PostgreSQL dans les bibliothèques du serveur d’application
	\end{itemize}
\end{frame}

\section{Architecture}
\begin{frame}
	\frametitle{Deux patterns principaux}
	Deux patterns courants liés à persistance : \emph{Active Record} et \emph{DAO}
	\begin{itemize}
		\item Active Record : chaque objet responsable de sa propre persistance
		\item Chaque objet contient méthodes CRUD : Create, Read, Update, Delete
		\item DAO ? \pause Data Access Object \pause
		\item Classes dédiées à interaction avec BD
		\item Permet isolation de cet aspect
		\item Typiquement utilisé avec Data Transfer Object
	\end{itemize}
	\begin{block}{Références}
		\begin{itemize}
			\item \href{http://gen.lib.rus.ec/book/index.php?md5=37E9F4F25E3C5609E14415472408B80C}{Patterns of Enterprise Application Architecture}
			\item \href{https://www.oracle.com/technetwork/java/dataaccessobject-138824.html}{Core J2EE Patterns - Data Access Object}
		\end{itemize}
	\end{block}
\end{frame}

\appendix

\section{Licence}
\begin{frame}
	\frametitle{Licence}
	Cette présentation, et le code LaTeX associé, sont sous \href{https://opensource.org/licenses/MIT}{licence MIT}. Vous êtes libres de réutiliser des éléments de cette présentation, sous réserve de citer l’auteur.
	
	Le travail réutilisé est à attribuer à \href{http://www.lamsade.dauphine.fr/~ocailloux/}{Olivier Cailloux}, Université Paris-Dauphine.
\end{frame}
\end{document}

\begin{frame}
	\frametitle{}
	\begin{itemize}
		\item 
	\end{itemize}
	\begin{block}{}
		
	\end{block}
\end{frame}

\section{Bibliographie}
\begin{frame}[allowframebreaks]
	\frametitle{Bibliographie}
	\def\newblock{\hskip .11em plus .33em minus .07em}
% 	\bibliography{zotero}
\end{frame}

\section{Autres}
\begin{frame}
	\frametitle{}
	\begin{itemize}
		\item 
	\end{itemize}
	\begin{block}{}
		
	\end{block}
\end{frame}
\end{document}
